\section{Discussion}

Ultra-short PPG ở cả Awake và Drowsy chứa tổ chức động lực học không được PPS null đã kiểm tra tái tạo đầy đủ. Drowsiness liên quan đến những thay đổi phối hợp về khả năng dự báo ở horizon hữu hạn, tổ chức tái diễn và mức phân kỳ quỹ đạo cục bộ ước lượng. Trong các kiểm tra sensitivity, hướng được duy trì tốt hơn độ lớn và mức hỗ trợ thống kê. Các thuộc tính bổ sung này cần được phân biệt theo estimand và thang động lực học, phù hợp với literature về predictability và nonlinearity~\cite{porta2007_prediction,porta2015,calderonjuarez2023,hu2009}.

\subsection{Changes in PPG Dynamics Across the Awake--Drowsy Transition}

Mean CC giảm và Mean NRMSE tăng cho thấy diễn tiến tương lai của PPG ở Drowsy được dự báo kém hơn từ các trạng thái lân cận, tương ứng với finite-horizon forecastability thấp hơn. Literature về local prediction của chuỗi khoảng nhịp phân biệt đánh giá khả năng dự báo với kiểm tra nonlinearity bằng surrogate~\cite{porta2007_prediction}. Faes et al.\ phân tách predictability trong mạng brain--heart thành các thành phần liên quan đến lịch sử của chính quá trình và thông tin từ các quá trình khác~\cite{faes2016}. Điều này đặt forecastability trong bối cảnh tổ chức theo thời gian của những hệ tương tác. Phân tích PPG đơn biến hiện tại chưa phân giải được thay đổi đến từ động lực học nội tại, tương tác giữa các hệ hay đóng góp tương đối của từng nguồn.

DET giảm tại embedding nominal phù hợp với tổ chức tái diễn theo đường chéo thấp hơn ở Drowsy; $L_{\mathrm{mean}}$ cùng hướng nhưng có mức hỗ trợ thống kê yếu hơn sau BH-FDR. Các nghiên cứu RQA trong sleep cho thấy recurrence organization thay đổi theo trạng thái và phụ thuộc vào cách tái dựng~\cite{rolink2015,martingonzalez2018}. Sự tái diễn theo đường chéo và sự lưu lại cục bộ theo chiều thẳng đứng có thể thay đổi khác nhau~\cite{martingonzalez2018}, nên xu hướng LAM/TT tăng không mâu thuẫn với DET giảm. Các thay đổi laminar vẫn là secondary do thiếu hỗ trợ sau BH-FDR. DET cần được hiểu theo recurrence geometry, không phải physical determinism; LAM/TT cũng không tự xác định nonlinearity. Literature sử dụng DET/LAM cùng surrogate testing củng cố sự phân biệt giữa giá trị descriptor và bằng chứng đối với một null cụ thể~\cite{calderonjuarez2023}.

LLE thấp hơn ở Drowsy biểu thị mức local trajectory divergence ước lượng thấp hơn trong fit interval đã khảo sát. Ước lượng này phụ thuộc vào thang quan sát, embedding, chiều dài dữ liệu, preprocessing và đặc tính tín hiệu; các nghiên cứu trên chuỗi nhịp tim cũng cho thấy cần xét thành phần tần số và thang động lực học khi diễn giải Lyapunov-related measures~\cite{yeragani2004,hu2009}. Forecastability giảm và LLE giảm không mâu thuẫn, vì Simplex và Rosenstein LLE ước lượng hai thuộc tính khác nhau. Simplex đánh giá mức độ các trạng thái lân cận cung cấp diễn tiến tương lai phù hợp với quan sát qua những horizon hữu hạn; LLE tóm tắt tốc độ phân kỳ cục bộ trong một khoảng fit cụ thể. Phân kỳ chậm hơn trong khoảng đó không bảo đảm dự báo chính xác hơn trên toàn bộ tập horizons. Hai estimands khai thác những khía cạnh khác nhau của thời gian và hình học quỹ đạo, phù hợp với literature về nonlinear dynamics ở nhiều thang~\cite{porta2015,hu2009}. Cùng với DET tại nominal embedding, chúng mô tả những chiều bổ sung của sự tái tổ chức PPG.

\subsection{Robustness and Finite-Window Effects}

Cửa sổ ngắn giúp định vị quá trình suy giảm vigilance theo thời gian, nhưng giới hạn số trạng thái, recurrence structures và trajectory pairs dùng cho ước lượng. Các hướng Mean CC giảm, Mean NRMSE tăng, DET giảm và LLE giảm được duy trì ở $30$, $60$, $120$ và $180$~s; độ bất định lớn hơn ở $30$~s và nhìn chung thấp hơn ở các độ dài còn lại. Tuy nhiên, mức giảm bất định không đồng đều giữa các metric, với DET tại $180$~s vẫn có khoảng bất định tương đối rộng. Sviridova et al.\ cho thấy độ dài PPG cần thiết để kiểm soát sai số RQA phụ thuộc vào metric và mức nhiễu, với DET cần đoạn ngắn hơn các diagonal-length measures trong điều kiện khảo sát~\cite{sviridova2022}. Điều này hỗ trợ đánh giá finite-window reliability theo từng descriptor. $60$~s là lựa chọn cân bằng thực tế giữa định vị theo thời gian và độ bất định trong dataset này, không phải độ dài tối ưu phổ quát.

Các label rules thay thế giữ cả bốn hướng khi loại vùng gần chuyển tiếp hoặc chỉ giữ các episode đủ dài, gợi ý pattern không chỉ xuất hiện dưới định nghĩa nhãn primary. Những kiểm tra này không thay thế việc xác nhận nhãn bằng phép đo khách quan. Bootstrap CI của trung vị chênh lệch paired và signed-rank inference dựa trên thứ hạng không kiểm tra cùng một đại lượng, nên CI chứa $0$ có thể đi cùng $q_{\mathrm{BH}}<0.05$. Với reconstruction, Mean CC, Mean NRMSE và LLE giữ hướng trong grid đã kiểm tra, nhưng điều này không bảo đảm độ lớn hoặc hỗ trợ thống kê ngoài grid. Literature về RQA và Lyapunov-related measures cũng nhấn mạnh sự phụ thuộc vào reconstruction, recurrence threshold, thành phần tín hiệu và thang phân tích~\cite{martingonzalez2018,yeragani2004,hu2009}. DET nhạy rõ với embedding delay: hiệu ứng âm tại nominal delay tiến gần $0$ và đổi dấu ở các delay lân cận, nên finding này cần gắn với cấu hình nominal.

Sensitivity đối với unknown repeated-session dependence giữ hướng của cả bốn headline effects dưới hypothetical clustering đã khảo sát, trong khi độ lớn thay đổi và hỗ trợ thống kê suy giảm. Hướng được duy trì tốt hơn inferential support, nhưng suy luận cấp participant vẫn chưa khả dụng.

\subsection{Physiological Interpretation of the Awake--Drowsy Transition}

Pattern tổng hợp phù hợp với sự tái cấu hình của PPG dynamics trong bối cảnh điều hòa tim mạch đang chuyển tiếp, thay vì một thay đổi đơn điệu theo complexity. Cách diễn giải này dựa trên bối cảnh wake-to-sleep transition, khi điều hòa trung ương, autonomic và tim mạch có thể thay đổi không đồng bộ~\cite{dezambotti2018,tobaldini2013}. Thay đổi tim mạch có thể bắt đầu trước sleep onset, bị gián đoạn bởi repeated arousals và tiếp tục khi giấc ngủ ổn định hơn~\cite{shinar2006,carrington2005}. Tsai et al.\ ghi nhận nhịp tim giảm trước cortical N1 onset ở người ngủ tốt, nhưng chậm hơn ở nhóm sleep-onset insomnia~\cite{tsai2019}, cho thấy cardiovascular de-arousal không nhất thiết trùng với cortical staging. Literature này cung cấp bối cảnh cho thay đổi phối hợp của các descriptor, chưa xác định cơ chế tạo ra chúng.

Symbolic dynamics cung cấp bối cảnh hỗ trợ: $0V$ tăng, $2V$ giảm, còn $1V$ gần $0$ và kém nhất quán hơn. Porta et al.\ ghi nhận $0V$ tăng và $2UV$ giảm trong chuỗi khoảng nhịp khi graded head-up tilt~\cite{porta2007_symbolic}. Sự tương đồng về hướng chỉ là analogy do khác điều kiện sinh lý và định nghĩa nhóm symbolic. Literature về sleep onset lại ghi nhận xu hướng thường được diễn giải là tăng vagal contribution hoặc giảm LF/HF~\cite{shinar2006,carrington2003}. Cùng với diễn tiến de-arousal không đồng bộ~\cite{tsai2019}, những quan sát này cho thấy symbolic pattern chưa xác định một ANS shift đơn giản. Hỗ trợ thống kê chưa nhất quán cũng giới hạn vai trò của symbolic ở lớp supportive.

Các descriptor của ultra-short PPG có thể bổ sung cho HR, HRV và hình thái waveform. Vaussenat et al.\ cũng khảo sát các đạo hàm HRV kết hợp với nhóm đặc trưng khác~\cite{vaussenat2026}. Đây là hướng nghiên cứu giá trị bổ sung của thông tin theo thời gian, cần được kiểm tra riêng cho PPG. Bài này chưa xây classifier hoặc đánh giá real-time/early detection; hàm ý ứng dụng nằm ở nghiên cứu nonlinear features và validation tiếp theo.

\subsection{Limitations of the Present Study}

Cohort nhỏ và chỉ gồm người trẻ khỏe mạnh, nên khả năng khái quát sang độ tuổi, tình trạng sức khỏe, sensor hoặc protocol khác còn hạn chế. Nhãn Awake/Drowsy dựa trên video và KSS chủ quan, không có PSG; Drowsy vì vậy không đồng nhất với N1 hoặc NREM. PPG là phép đo tim mạch ngoại vi, không trực tiếp đo hoạt tính trung ương hoặc autonomic, và symbolic metrics không phải phép đo trực tiếp sympathetic/parasympathetic activity.

Preprocessing, chất lượng tín hiệu, window length, phase-space reconstruction, recurrence threshold, định nghĩa neighborhood và LLE fit interval giới hạn các ước lượng trên dữ liệu hữu hạn~\cite{martingonzalez2018,sviridova2022,yeragani2004,hu2009}. Cố định tham số và QC không loại bỏ estimator-specific dependence hoặc residual nonstationarity, và sensitivity chỉ bao phủ các cấu hình đã khảo sát. Bác bỏ PPS là bằng chứng chống lại noisy periodic-orbit-plus-uncorrelated-noise null đã kiểm tra, không chứng minh deterministic chaos và không loại trừ pseudoperiodicity phức tạp hơn có cấu trúc tương quan.

$20$ recording sessions đến từ $10$ participants, nhưng thiếu liên kết participant--session nên primary inference vẫn ở cấp session. Cấu trúc phụ thuộc giữa repeated sessions vẫn chưa biết. Hypothetical clustering chỉ khảo sát sensitivity, không khôi phục participant identity hoặc chứng minh session independence; suy luận cấp participant chưa khả dụng. Nghiên cứu chưa có simultaneous EEG, respiration, continuous blood pressure hoặc direct autonomic measures, và chưa tái lập trên independent dataset. Các findings là associations, không xác lập causal physiological mechanisms. Multimodal validation cùng dữ liệu có participant--session linkage cần thiết để kiểm chứng diễn giải sinh lý, phân biệt thay đổi theo trạng thái với khác biệt cá thể và đánh giá khả năng khái quát.